% ===========================================================================
% DA Galois Connection Theorem
% ===========================================================================
\subsection{DA as a Canonical Instance of Abstract Interpretation}
\label{sec:galois}

The relationship between NeuroPLC and abstract interpretation
\cite{cousot1977abstract} extends beyond methodological inspiration.
We prove that Doubleton Arithmetic constitutes a \textbf{canonical
instance} of the Cousot abstract interpretation framework---with a
defined Galois connection $(\alpha, \gamma)$ that establishes DA as
the \textit{optimal} abstract domain for the $C^2$ function class
that KAN activations belong to.

\subsubsection{The DA Galois Connection}

\begin{definition}[DA Abstract Interpretation]
\label{def:da_galois}
The DA abstract interpretation is the quadruple $(D_C, D_A, \alpha,
\gamma)$ where:

\begin{itemize}
    \item \textbf{Concrete domain} $D_C$: the set of $C^2$ functions
    $f: [x_0-r, x_0+r] \to \mathbb{R}$ with
    $M_2(f) = \sup_x |f''(x)| < \infty$, ordered by the pointwise
    order $f_1 \leq_C f_2 \iff \forall x: f_1(x) \leq f_2(x)$.

    \item \textbf{Abstract domain} $D_A$: the set of doubleton
    pairs $(c, R) \in \mathbb{R} \times \mathbb{R}_{\geq 0}$,
    ordered by precision: $(c_1, R_1) \leq_A (c_2, R_2) \iff
    |c_1 - c_2| + R_1 \leq R_2$. Intuitively, $d_1 \leq_A d_2$
    means $d_2$ is a looser abstraction than $d_1$.

    \item \textbf{Abstraction function} $\alpha: D_C \to D_A$:
    \begin{equation}
    \alpha(f) = \bigl(f(x_0),\; |f'(x_0)| \cdot r + M_2(f) \cdot r^2/2\bigr)
    \label{eq:alpha_da}
    \end{equation}
    This maps a concrete function to the sound doubleton that
    over-approximates its output range on the given interval; the
    radius combines the first-order contribution $|f'(x_0)| \cdot r$
    with the second-order curvature term $M_2(f) \cdot r^2/2$ from
    the second-order Taylor bound.

    \item \textbf{Concretization function} $\gamma: D_A \to
    \mathcal{P}(D_C)$:
    \begin{equation}
    \gamma(c, R) = \bigl\{f \in D_C \mid |f(x_0) - c| + |f'(x_0)| \cdot r
    + M_2(f) \cdot r^2/2 \leq R\bigr\}
    \label{eq:gamma_da}
    \end{equation}
    This maps an abstract doubleton to the set of all concrete
    functions whose first-order-plus-curvature envelope fits within
    $R$ of the center $c$; it is the right adjoint generated by
    $\alpha$ (the coarsest concretization for which the Galois
    condition below holds).
\end{itemize}
\end{definition}

\vspace{4pt}
\begin{kingthm}[DA Galois Connection]
\label{thm:galois}
The pair $(\alpha, \gamma)$ defined above forms a \textbf{Galois
connection} between $D_C$ and $D_A$: $\forall f \in D_C,\;
\forall d \in D_A$:
\begin{equation}
\alpha(f) \leq_A d \quad\Longleftrightarrow\quad f \in \gamma(d)
\label{eq:galois_condition}
\end{equation}

Moreover, the DA abstract transformers are \textbf{Galois-monotone}:
for any DA operation $\texttt{op}_C: D_C^n \to D_C$, its abstract
counterpart $\texttt{op}_A = \alpha \circ \texttt{op}_C \circ
\gamma^n$ satisfies $\texttt{op}_A: D_A^n \to D_A$ and is monotone
with respect to $\leq_A$.
\end{kingthm}

\vspace{2pt}
\noindent\textit{Proof sketch.}
Four properties of the Galois connection are verified:

\begin{enumerate}[label=(\roman*)]
    \item \textbf{Soundness (extensiveness of $\gamma \circ \alpha$):}
    $\forall f \in D_C: f \in \gamma(\alpha(f))$. The second-order
    Taylor bound gives
    $|f(x) - f(x_0) - f'(x_0)(x-x_0)| \leq M_2(f)|x-x_0|^2/2$ for all
    $x \in [x_0-r, x_0+r]$, hence by the triangle inequality
    $|f(x) - f(x_0)| \leq |f'(x_0)| \cdot r + M_2(f) \cdot r^2/2$.
    Therefore $f$ satisfies the envelope condition with
    $c = f(x_0)$ and $R = |f'(x_0)|r + M_2(f)r^2/2$, i.e.
    $f \in \gamma(\alpha(f))$.

    \item \textbf{Adjunction (Galois condition):}
    $\forall f \in D_C,\; \forall d = (c,R) \in D_A$:
    $\alpha(f) \leq_A d \iff f \in \gamma(d)$.
    ($\Rightarrow$) Unrolling $\alpha(f) \leq_A (c,R)$ gives
    $|f(x_0) - c| + |f'(x_0)|r + M_2(f)r^2/2 \leq R$, which is
    exactly the $\gamma$ membership condition.
    ($\Leftarrow$) $f \in \gamma(c,R)$ is the same inequality, which
    is precisely $\alpha(f) \leq_A (c,R)$ by the definition of
    $\leq_A$ (Eq.~\ref{eq:alpha_da}). Hence the Galois condition
    holds by construction of $\gamma$ as the right adjoint of
    $\alpha$.

    \item \textbf{Optimality (tightness of $\alpha \circ \gamma$):}
    For every $d = (c,R) \in D_A$, the $\leq_A$-infimum of
    $\alpha(\gamma(d))$ is $d$ itself. Lower bound: by (ii),
    every $f \in \gamma(c,R)$ satisfies $\alpha(f) \leq_A (c,R)$.
    Attainment: the linear function
    $f^*(x) = c + (R/r)(x-x_0)$ satisfies
    $|f^*(x_0) - c| + |(f^*)'(x_0)| \cdot r + M_2(f^*) \cdot r^2/2
    = 0 + R + 0 = R$,
    so $f^* \in \gamma(c,R)$ and $\alpha(f^*) = (c, R)$. Hence no
    strictly tighter doubleton soundly abstracts every function in
    $\gamma(d)$; the abstraction is tight (this is the Galois
    connection form of the DA optimality proved in
    Corollary~\ref{cor:da-optimality}; note the $\alpha \circ \gamma
    = \text{id}$ \emph{insertion} form does not hold, since constant
    functions inside $\gamma(c,R)$ abstract to $(c,0)$).

    \item \textbf{Galois-monotonicity of DA operations:}
    \begin{itemize}
        \item \textbf{DA Addition}: $+(c_1,R_1)(c_2,R_2) = (c_1+c_2,
        R_1+R_2)$. If $(c_1,R_1) \leq_A (c'_1,R'_1)$ and
        $(c_2,R_2) \leq_A (c'_2,R'_2)$, then $|(c_1+c_2)-(c'_1+c'_2)|
        + (R_1+R_2) \leq (R'_1+R'_2)$ by triangle inequality.
        Monotonicity holds.

        \item \textbf{DA Scalar Multiply}: $k \cdot (c, R) = (k \cdot
        c, |k| \cdot R)$. Monotonicity follows from $|k| \geq 0$.

        \item \textbf{DA Element-wise Application}: DA($\phi$,
        $(c,R)$) $= (\phi(c),\; M_2(\phi) \cdot R^2 + |\phi'(c)|
        \cdot R)$. Monotonicity holds because increasing $c$ or $R$
        increases both terms. Verified by the coefficient analysis
        of Corollary~\ref{cor:da-optimality}.
    \end{itemize}
\end{enumerate}

The Galois connection holds \textbf{unconditionally} for DA on the
$C^2$ function class but composes across network layers \textbf{only}
when the architecture satisfies SVNN Conditions~1--2. This is exactly
the structural requirement that enables the per-layer Galois
connection to compose into the end-to-end guarantee of
Theorem~\ref{thm:svnn_conditions}. $\square$

\vspace{4pt}
\noindent\textbf{Positioning in Abstract Interpretation Theory.}
Theorem~\ref{thm:galois} grounds NeuroPLC's compiler in the standard
formal framework of abstract interpretation. The DA abstract domain
$(\alpha, \gamma)$ is:

\begin{itemize}
    \item \textbf{The first Galois connection defined for neural
    network activation functions}---prior abstract-interpretation
    tools for neural networks (DeepPoly~\cite{singh2019abstract},
    CROWN~\cite{zhang2018crown}, AI\textsuperscript{2}%
    \cite{gowal2018scalable}) use abstract domains (intervals,
    zonotopes, polyhedra) but define these domains for \textit{ReLU}
    activations specifically, without the general $C^2$ Galois
    connection.

    \item \textbf{Optimal for its input class}---DA is the
    tightest possible abstract domain for $C^2$ univariate
    functions (Corollary~\ref{cor:da-optimality} proves coefficient
    optimality; the tightness of the de~Boor bound
    (Theorem~\ref{thm:tightness}) proves bound attainment).

    \item \textbf{The basis for a compositional certification
    theory}---SVNN Conditions~1--2 are precisely the architectural
    restrictions that make the per-operation Galois connections
    composable across layers (Theorem~\ref{thm:svnn_conditions},
    Step~3). This is a novel insight: the Cousot framework
    naturally explains \textit{why} operation separation enables
    compositional certification.
\end{itemize}

The Galois connection provides the formal justification for a claim
made earlier in the paper: NeuroPLC's compiler correctness is not
merely an engineering achievement but an instance of a \textit{general
abstract interpretation theory for neural architectures}. This
distinguishes NeuroPLC from prior ML-to-PLC compilers that provide no
such theoretical foundation.
